\chapter{Editorial and Factual Review}
\label{ch:audit}

\begin{chapterguide}
This appendix records the final review standard applied to the manuscript. It
distinguishes factual verification, source verification, argument review, and
visual review. It also states the remaining limits of a generated scholarly
manuscript so that readers can assess it responsibly.
\end{chapterguide}

A publication-ready philosophy book needs more than fluent prose. It needs
claims that can be traced, references that exist, arguments that do not
silently change level, and pages that can be read without visual distraction.
This appendix documents the review categories used for this edition.

\section{Factual and source review}

The historical claims in the manuscript were checked against primary
publications where available and against respected academic reference sources
for broader philosophical positions. The most consequential empirical cases
use the following source classes:

\begin{table}[ht]
\centering
\smalltable
\begin{tabularx}{\textwidth}{p{0.25\textwidth}Y Y}
\toprule
Area & Claim type audited & Source anchor\\
\midrule
Philosophy of science & Dates, positions, and scope of major arguments & Original books and articles; Stanford Encyclopedia entries\\
Physics & Brownian motion, electron, ether experiment & Einstein, Perrin, Thomson, Michelson and Morley; later historical scholarship\\
Biology & 1953 DNA structure paper and evidential role of diffraction & Watson and Crick (1953); associated primary papers and historical records\\
Medicine & 1948 controlled streptomycin investigation & Medical Research Council (1948), British Medical Journal\\
Psychology & 2015 replication project design and reported interpretation & Open Science Collaboration (2015), Science\\
Climate & AR6 synthesis and calibrated attribution language & IPCC (2023) Synthesis Report\\
Economics & Policy-regime change and Lucas critique & Lucas (1976)\\
Social science & Audit-study design and scope & Bertrand and Mullainathan (2004)\\
\bottomrule
\end{tabularx}
\caption{Source anchors for the factual case studies.}
\label{tab:audit-sources}
\end{table}

The manuscript avoids invented quotations and page numbers. Where a claim is
an interpretation of a historical episode, the text marks it as a
philosophical reading rather than presenting it as a direct quotation from the
source. Publication details should still be checked against the particular
edition used for a print edition or formal scholarly citation.

\section{Argument review}

The following checks were applied to the argument:

\begin{enumerate}
\item Major traditions are presented in their strongest recognizable form
before criticism.
\item Descriptive history, established scientific evidence, interpretation,
the author's original framework, and speculation are separated.
\item The framework has an explicit ontology, epistemology, axioms, inferential
rules, criteria, objections, and limits.
\item Formal notation is defined before use. Equations are used to clarify
Bayesian updating, measurement, growth, and causal inference rather than to
create a false appearance of certainty.
\item Causal, predictive, explanatory, pragmatic, and normative claims are not
treated as interchangeable.
\item Social and feminist approaches are included as contributions to
objectivity, not as decorative additions or caricatures.
\item Historical examples are used as stress tests, not as simplistic proofs
of a contemporary theory.
\item Claims are scoped where transport, intervention, or measurement is
limited.
\end{enumerate}

\section{Reference review}

The bibliography was checked for the existence of the cited works, author
names, publication years, titles, journals or publishers, and persistent
identifiers where included. The references are deliberately conservative:
when edition-specific page information was uncertain, page numbers were not
invented. Readers preparing a journal article, dissertation, or commercial
edition should verify the exact edition and citation style required by their
publisher.

A bibliography is not evidence by itself. A cited source can be interpreted
badly. The manuscript therefore uses citations as anchors for claims and
states when an argument is the author's own.

\section{Language and editorial review}

The manuscript was reviewed for:

\begin{itemize}
\item consistent use of \enquote{science}, \enquote{scientific},
\enquote{theory}, \enquote{model}, \enquote{evidence}, \enquote{truth}, and
\enquote{objectivity};
\item avoidance of unexplained specialist terms where a first-principles
definition is needed;
\item separation of paragraphs that make different argumentative moves;
\item consistent terminology for Reflexive Constraint Realism;
\item removal of unsupported superlatives and claims of finality;
\item clear distinction between a method's scope and its philosophical
interpretation;
\item discussion questions that follow from the chapter rather than repeat its
headings.
\end{itemize}

The book intentionally uses a restrained tone. It argues for a position without
claiming that the position has settled every debate.

\section{Visual and typesetting review}

The LaTeX project uses a 6 by 9 inch book page, restrained colour, readable
margins, consistent headers and footers, and tables with fixed column logic.
The cover uses a solid background and only the title, subtitle, and author
name. No decorative lines, overlapping shapes, or crowded geometry are used.

The final PDF should be rendered page by page and checked for:

\begin{enumerate}
\item clipped or overflowing text;
\item table collisions or unreadable cells;
\item broken equations or missing glyphs;
\item bad line breaks in headings and captions;
\item inconsistent page numbering and running headers;
\item blank pages that do not serve a chapter transition;
\item index and table-of-contents links;
\item cover spacing and contrast.
\end{enumerate}

This appendix is a review record, not a guarantee that no future copy-editing
decision is possible. Scholarly publication includes expert review, source
checking by specialists in each case domain, and updates as evidence changes.
The purpose of this edition is to provide a coherent, transparent, editable
manuscript that makes those next checks possible.

\section*{Review conclusion}

The manuscript meets its stated design target when it is compiled with the
provided source and the rendered PDF passes the visual checks above. Its
central philosophical proposal is original in formulation and argumentative
organization, while its historical materials are grounded in established
literature. The remaining limitations are the normal limits of a single-author
work spanning many sciences: domain specialists may refine examples, later
research may revise empirical assessments, and the framework itself remains
open to criticism.

